\section{Three different meanings of coordinates}
\label{sec:coordinate-layers}

The fourth question of this paper is not where to operate, how to move, or
which machine to build, but \emph{in which variables the other three questions
should be asked}. A useful first safeguard is to separate three constructions
that are often all called a frame.

\begin{description}
 \item[Spatial electromagnetic frame.] Clarke and Park transformations express
 the same two-dimensional stator vector in stationary \(\alpha\beta\),
 rotor-oriented \(dq\), or flux-oriented \(ft\) basis vectors. Voltage,
 current, and stator flux all rotate together.
 \item[State or task coordinates.] A map from
 \((i_{\mathrm d},i_{\mathrm q},i_{\mathrm{exc}})\) to torque, flux, and
 allocation variables reshapes the three-dimensional current space. It may
 mix a rotor current with stator components and is not a physical rotation of
 the air-gap field.
 \item[Virtual control coordinates.] A state-dependent input map may replace
 physical converter voltages by commands that request torque rate, flux rate,
 and torque-neutral motion. This is a change of input channels, not a change
 of spatial basis.
\end{description}

\begin{figure}[tb]
\centering
\begin{tikzpicture}[node distance=8mm and 8mm,>=Latex,font=\small]
\node[draw,rounded corners,fill=MidnightBlue!8,align=center,minimum width=3.4cm,minimum height=1.15cm] (space)
 {spatial frame\\\(abc\to\alpha\beta\to dq\to ft\)};
\node[draw,rounded corners,fill=ForestGreen!9,align=center,right=of space,minimum width=3.4cm,minimum height=1.15cm] (task)
 {state/task chart\\\(\vect i\to(M,\psi,\sigma)\)};
\node[draw,rounded corners,fill=Orange!12,align=center,right=of task,minimum width=3.4cm,minimum height=1.15cm] (input)
 {virtual inputs\\\(\vect u\to(v_M,v_\psi,v_N)\)};
\draw[->,thick] (space) -- node[above]{measured state} (task);
\draw[->,thick] (task) -- node[above]{Jacobian} (input);
\node[below=8mm of space,align=center,text width=3.2cm]{same 2-D electromagnetic vector\\different rotating basis};
\node[below=8mm of task,align=center,text width=3.2cm]{different labels for the\\3-D current manifold};
\node[below=8mm of input,align=center,text width=3.2cm]{different directions of\\converter action};
\end{tikzpicture}
\caption{Three meanings of coordinates. Only the first column is a spatial
rotation of the stator field.}
\label{fig:coordinate-layers}
\end{figure}


Confusing these levels hides costs. A Park transform needs an angle but is
linear once that angle is known. A task chart may require a nonlinear inverse.
A virtual-input map may require a matrix inversion at every sample. The
machine decides which representation is useful; no coordinate name removes
the underlying sensing, inversion, or actuation burden.

\subsection{Why the rotor frame remains unusually effective}

For a balanced sinusoidal space vector
\[
 \vect x_{\alpha\beta}(t)=X
 \begin{bmatrix}\cos\theta_{\mathrm e}(t)&\sin\theta_{\mathrm e}(t)\end{bmatrix}^{\mathsf T},
\]
the power-invariant Clarke map first removes the redundant zero-sequence
direction and the Park map gives
\[
 \vect x_{dq}=R(-\gls{sym:thetae})\vect x_{\alpha\beta}
 =\begin{bmatrix}X&0\end{bmatrix}^{\mathsf T}.
\]
Sinusoids become constants, rotor saliency axes become fixed, and the
inductance matrix is stationary. Two constant references can therefore be
tracked by ordinary PI current loops. This is why \gls{dq} solves a
\emph{representation problem} extremely well. It does not follow that
\((i_{\mathrm d},i_{\mathrm q})\) are also the best optimization outputs or
that three EESM currents should be controlled as three equally important
channels.

\begin{figure}[tb]
\centering
\begin{tikzpicture}[>=Latex,font=\small]
\node[draw,circle,minimum size=10mm] (a) at (0,0) {\(a\)};
\node[draw,circle,minimum size=10mm] (b) at (0,-1.2) {\(b\)};
\node[draw,circle,minimum size=10mm] (c) at (0,-2.4) {\(c\)};
\node[draw,rounded corners,fill=MidnightBlue!8,minimum width=2.4cm,minimum height=1cm] (cl) at (3,-1.2) {Clarke};
\node[draw,rounded corners,fill=ForestGreen!9,minimum width=2.4cm,minimum height=1cm] (pa) at (6.3,-1.2) {Park \(R(-\theta_e)\)};
\node[draw,rounded corners,fill=Orange!12,minimum width=2.8cm,minimum height=1cm] (dc) at (9.8,-1.2) {\(x_d=X,\ x_q=0\)};
\foreach \n in {a,b,c}\draw[->] (\n.east)--(cl.west);
\draw[->,thick] (cl)--node[above]{\(\alpha,\beta\) sinusoidal}(pa);
\draw[->,thick] (pa)--node[above]{co-rotate}(dc);
\draw[->,Purple,very thick] (3,-3.0)--++(1.2,0.55) node[right]{\(\vect x_{\alpha\beta}\)};
\draw[->,gray,thick] (3,-3.0)--++(1.2,0) node[right]{\(\alpha\)};
\draw[->,gray,thick] (3,-3.0)--++(0,1.0) node[above]{\(\beta\)};
\draw[->,Purple,very thick] (7.5,-3.0)--++(1.2,0) node[right]{\(d,\vect x\)};
\draw[->,gray,thick] (7.5,-3.0)--++(0,1.0) node[above]{\(q\)};
\end{tikzpicture}
\caption{Clarke removes the redundant three-phase direction; Park follows the
rotating sinusoid so the steady vector becomes constant.}
\label{fig:abc-clarke-park}
\end{figure}


\subsection{An arbitrary synchronous frame and its price}

Let the absolute frame angle be
\[
 \gls{sym:gammaframe}(t)=\gls{sym:thetae}(t)+\gls{sym:frameangle}(t),\qquad
 R(\delta)=\begin{bmatrix}\cos\delta&-\sin\delta\\
 \sin\delta&\cos\delta\end{bmatrix},\quad
 \gls{sym:jrot}=\begin{bmatrix}0&-1\\1&0\end{bmatrix}.
\]
The forward and inverse transformations are
\begin{equation}
 \vect x_\gamma=R(-\delta)\vect x_{dq},\qquad
 \vect x_{dq}=R(\delta)\vect x_\gamma.
 \label{eq:arbitrary-frame-map}
\end{equation}
Because \(\dot R(-\delta)=-\dot\delta\,\matr J R(-\delta)\),
\begin{equation}
 \boxed{\dot{\vect x}_\gamma
 =R(-\delta)\dot{\vect x}_{dq}
 -\dot\delta\,\matr J\vect x_\gamma.}
 \label{eq:arbitrary-frame-rate}
\end{equation}
The final term is not a machine force. It is the apparent motion caused by
the moving ruler. For \(\delta=\mathrm{const.}\) the frame is rotor-fixed and
the term vanishes. If \(\delta=\delta(\vect i,\boldsymbol\psi)\), then
\(\dot\delta=\nabla\delta^{\mathsf T}\dot{\vect x}\) and the frame is only
synchronous with the selected flux during steady operation. We therefore use
\emph{rotor-fixed}, \emph{rotor-synchronous}, and
\emph{stator-flux-oriented} as distinct terms throughout.

\begin{quote}\itshape
Rotating a basis may simplify one algebraic relation, but a time-varying basis
adds a connection term to the dynamics. There is no free coordinate system.
\end{quote}
